%% Рукопись для Journal of Healthcare Informatics Research (Springer).
%% Русская версия. Шаблон Springer Nature sn-jnl, стиль ссылок sn-basic
%% (нумерованные ссылки в квадратных скобках).
%%
%% Сборка: pdfLaTeX. Три строки ниже (fontenc/inputenc/babel) нужны только для
%% кириллицы, в английской версии их нет. Опция Numbered обязательна, иначе
%% sn-basic идет в режиме автор-год и natbib выдает ошибку.
%%
%% Держать manuscript_ru.md, manuscript_en.md и manuscript_en.tex
%% синхронно при правке содержания.

\documentclass[pdflatex,sn-basic,Numbered]{sn-jnl}

%%%% Стандартные пакеты шаблона
\usepackage{graphicx}%
\usepackage{multirow}%
\usepackage{amsmath,amssymb,amsfonts}%
\usepackage{amsthm}%
\usepackage[title]{appendix}%
\usepackage{xcolor}%
\usepackage{textcomp}%
\usepackage{manyfoot}%
\usepackage{booktabs}%

%%%% Только для русской версии
\usepackage[T2A]{fontenc}%
\usepackage[utf8]{inputenc}%
\usepackage[english,russian]{babel}%

\raggedbottom

\begin{document}

\title[Факторы рецидивирующего течения ДКА]{Клинические факторы рецидивирующего течения диабетического кетоацидоза: анализ российской когорты методами статистики и интерпретируемого машинного обучения}

\author*[1]{\fnm{Олег} \sur{Садыхов}}\email{ovelsad@niuitmo.ru}

\author[1]{\fnm{Василий} \sur{Леоненко}}\email{vnleonenko@itmo.ru}

\affil*[1]{\orgdiv{Факультет технологий искусственного интеллекта}, \orgname{Университет ИТМО}, \orgaddress{\street{Кронверкский пр., д. 49}, \city{Санкт-Петербург}, \postcode{197101}, \country{Россия}}}

\abstract{\textbf{Цель.} Выявить клинические факторы, связанные с рецидивирующим течением диабетического кетоацидоза (ДКА), на российской когорте и оценить методами интерпретируемого машинного обучения, насколько эти группы различаются по клиническому профилю при поступлении.

\textbf{Методы.} Ретроспективное одноцентровое исследование, предикторы и исход относятся к одной госпитализации, времени до события в данных нет, поэтому модель диагностическая, а не прогностическая на временном горизонте. Исход установлен у 273 пациентов из 289 (рецидивирующее течение 87, единичный эпизод 186). Группы сравнивали по фиксированному статистическому протоколу с поправкой Бенджамини-Хохберга. Независимо строили четыре интерпретируемые модели на 25 не зависящих от исхода признаках без мультиколлинеарности, разбиение выборки предшествовало всей предобработке. Вклад признаков оценивали методом SHAP и сопоставляли с результатами классической статистики. Заранее заданы две гипотезы, о связи приема алкоголя и о превосходстве сложных моделей над логистической регрессией.

\textbf{Результаты.} С исходом значимо связаны 10 признаков. Разделяющая способность умеренная, ROC-AUC 0.71-0.78 по вложенной кросс-валидации и 0.72-0.80 на отложенном тесте. В топ-10 по SHAP одновременно у всех четырех моделей входят пять признаков, четыре из них статистически значимы, пятый (глюкоза при поступлении) выявлен только многомерным подходом. При бутстрэпе общими для всех моделей стабильно остаются два признака из пяти (89-92\% повторов). Клиническая гипотеза подтвердилась, отношение шансов 2.27 (95\% ДИ 1.23-4.20). Методологическая гипотеза не подтвердилась, критерий DeLong не выявил значимых различий между моделями.

\textbf{Заключение.} Клинический профиль умеренно разделяет группы, а верхняя часть списков факторов, полученных двумя независимыми методологиями, совпадает.}

\keywords{диабетический кетоацидоз, рецидивирующее течение, факторы риска, интерпретируемое машинное обучение, SHAP, диагностическая модель}

\maketitle

\section{Введение}\label{sec1}

Диабетический кетоацидоз - жизнеугрожающее осложнение сахарного диабета, для которого характерны гипергликемия, метаболический ацидоз и повышенная концентрация кетоновых тел \cite{ref1}. За период с 2003 по 2014 год в США зарегистрировано 1 760 101 первичная госпитализация по поводу ДКА, а число выписок с основным диагнозом выросло со 118 808 в 2003 году до 188 965 в 2014 году \cite{ref2}. После купирования острого состояния единичные случаи могут перерастать в повторные эпизоды, а рецидивирующее течение сопряжено с существенно более высокой летальностью \cite{ref3,ref4}.

Повторные эпизоды ДКА редко объясняются одной причиной. К факторам риска повторных госпитализаций относят молодой возраст, психиатрические расстройства, инфекции, несоблюдение режима терапии, злоупотребление алкоголем и психоактивными веществами, а также социально-экономические обстоятельства \cite{ref5,ref6,ref7}. Набор факторов разнороден, результаты разных работ плохо согласуются между собой, а на российской популяции состав факторов рецидивирующего течения не описан. Несмотря на десятилетия клинического опыта, оптимальные стратегии ведения ДКА остаются предметом дискуссий \cite{ref8}.

Ближе всего к нашей задаче стоит наблюдательное исследование типа случай-контроль, где пациентов с рецидивирующим и единичным течением уравняли по возрасту, длительности диабета, полу и этнической принадлежности \cite{ref9}. После выравнивания заметных различий между группами почти не осталось, а основным поводом к эпизоду в обеих группах служил пропуск инъекций инсулина. При 40 парах пациентов такой результат трудно трактовать однозначно, он может означать и отсутствие различий, и нехватку мощности. Вопрос о том, какие характеристики пациента связаны с рецидивирующим течением, остается открытым.

Машинное обучение позволяет посмотреть на ту же задачу с другой стороны. Опубликованные модели по ДКА достигают ROC-AUC порядка 0.82-0.85 \cite{ref10,ref11}, но решают другую задачу, предсказывают первый эпизод или факт госпитализации в заданном временном окне, опираются на крупные зарубежные базы электронных медицинских карт с высокой долей пропусков, и на других выборках их метрики обычно не воспроизводятся \cite{ref11}. Кроме того, ранее показано \cite{ref12,ref13}, что многие модели работают как черный ящик, их выводы трудно интерпретировать, что снижает доверие врачей и становится барьером для применения. Современные работы в диабетологии применяют методы объяснения предсказаний, прежде всего SHAP и LIME \cite{ref14,ref15}, но для рецидивирующего течения ДКА такие инструменты почти не использовались.

В настоящей работе мы выявляем клинические факторы, связанные с рецидивирующим течением ДКА, на российской когорте и оцениваем, насколько по клиническому профилю пациента при поступлении можно отличить рецидивирующее течение от единичного эпизода. Ключевая особенность подхода в том, что состав факторов мы получаем двумя независимыми путями, классической статистикой по согласованному протоколу и методом SHAP поверх интерпретируемых моделей машинного обучения, после чего сопоставляем результаты. Такое сравнение показывает, в какой мере состав факторов зависит от выбранного подхода к анализу, что особенно важно при малой выборке.

Мы проверяем две заранее заданные гипотезы, первая из которых клиническая, касающаяся приема алкоголя, который в литературе описан как фактор рецидивирующего течения наряду с приемом психоактивных веществ, а избыток алкоголя назван одним из частых провоцирующих факторов эпизода \cite{ref7}. Связь подтверждают и работы по повторным госпитализациям \cite{ref9,ref16}. Остается выяснить, переносится ли она на российскую выборку. Методологическая гипотеза касается выбора модели, систематический обзор клинических прогностических моделей не обнаружил преимущества машинного обучения над логистической регрессией \cite{ref17}, однако на задаче рецидивирующего ДКА это не проверялось. Ответ имеет практическое следствие, если сложность не дает выигрыша, предпочтительна простая и интерпретируемая модель.

\section{Методы}\label{sec2}

\subsection{Дизайн и источник данных}\label{subsec2_1}

Ретроспективное одноцентровое исследование на обезличенных данных с поперечным установлением исхода, и предикторы, и исход относятся к одной госпитализации с ДКА, периода наблюдения между ними нет. В анализ включены пациенты, госпитализированные с диабетическим кетоацидозом с 2019 по 2025 год. Критерии включения, возраст 18 лет и старше и подтвержденный диагноз ДКА. Единственный критерий исключения - неустановленный исход. Исходная выборка составила 289 пациентов, исход установлен у 273.

\subsection{Исход}\label{subsec2_2}

Целевая переменная - рецидивирующее течение ДКА (0 - единичный эпизод, 1 - рецидивирующее течение). Пациента относили к группе рецидивирующего течения, если суммарно у него зарегистрировано более одного эпизода ДКА. Счет эпизодов складывался из числа прежних обращений пациента в эту больницу с тем же диагнозом по медицинской документации, данных анамнеза при поступлении, если ранее пациент здесь не наблюдался, и текущей госпитализации. Единичным считали первое и единственное обращение.

Исход устанавливали ретроспективно, поэтому на момент измерения предикторов он в большинстве случаев уже реализовался. Фиксированного окна наблюдения, дат отдельных эпизодов и времени до события в данных нет, цензурирование отсутствует. По классификации TRIPOD такая модель диагностическая, она распознает уже наступившее состояние, а не предсказывает будущее событие. Прямого сопоставления с моделями, предсказывающими госпитализацию за 30, 90 или 180 дней, она не допускает. Решаемая задача - разделение рецидивирующего и единичного течения по клиническому профилю пациента и выявление связанных с ним факторов. Следствия для интерпретации разобраны в разделе ограничений.

\subsection{Гипотезы}\label{subsec2_3}

Обе гипотезы сформулированы заранее, до анализа данных, на основании литературы.

\textbf{Клиническая гипотеза.} H0: прием алкоголя за сутки до эпизода не связан с рецидивирующим течением (отношение шансов равно 1). H1: связь есть (отношение шансов отличается от 1). Проверка - критерий для четырехпольной таблицы по протоколу, мера эффекта - отношение шансов с 95\% доверительным интервалом, с поправкой Бенджамини-Хохберга на множественные сравнения.

\textbf{Методологическая гипотеза.} H0: сложные модели (случайный лес, градиентный бустинг) и логистическая регрессия не различаются по разделяющей способности. H1: сложные модели превосходят логистическую регрессию по ROC-AUC. Проверка - парный критерий DeLong для связанных ROC-кривых на одних и тех же пациентах, отдельно на out-of-fold обучающей части и на отложенном тесте, с поправкой Бенджамини-Хохберга.

Других гипотез мы не выдвигали. Все прочие статистические сравнения носят описательный характер, а признаки со значимыми различиями между группами трактуются как результат разведочного анализа, требующий подтверждения на независимой выборке.

\subsection{Предикторы}\label{subsec2_4}

Использованы показатели, собранные при госпитализации, демографические данные (возраст, пол), анамнез и терапия сахарного диабета (тип, длительность, возраст манифестации, вид инсулинотерапии, суточная доза), лабораторные показатели при поступлении (газы крови, электролиты, креатинин, мочевина, глюкоза, HbA1c, липидограмма), осложнения (хроническая болезнь почек по стадии и альбуминурии, диабетические ретинопатия и полинейропатия) и фактор образа жизни (прием алкоголя).

Модели обучали на одном наборе признаков - 25 показателей без выраженной мультиколлинеарности. Возраст почти линейно связан с возрастом манифестации и длительностью диабета (коэффициент корреляции Пирсона около 0.99), поэтому возраст манифестации исключен как линейно зависимый. Состав набора определялся только взаимной коррелированностью предикторов и от исхода не зависел, поэтому отбор признаков не вносит смещения в оценку качества.

\subsection{Размер выборки}\label{subsec2_5}

Размер выборки не рассчитывали заранее, исследование ретроспективное, поэтому в анализ вошли все доступные наблюдения - 273 пациента с установленным исходом (доля рецидивирующего течения 31.9\%, 87 событий). При 87 событиях на 25 признаков приходится мало наблюдений редкого класса, что повышает риск переобучения, поэтому мы ограничили размерность, применяли регуляризацию и взвешивание классов. Признаков переобучения искали тремя способами, сопоставлением оценок на вложенной кросс-валидации и отложенном тесте, out-of-bag оценкой случайного леса и 95\% доверительными интервалами метрик по бутстрэпу.

\subsection{Пропущенные данные}\label{subsec2_6}

Доля пропусков по показателям менялась от единиц процентов до 58\% (HbA1c, натрий, калий, лактат). Мы исходили из допущения, что пропуски случайны при условии наблюдаемых переменных, и учитывали их панельную структуру, когда лабораторные показатели пропадают группами. Сравнили стратегии импутации, медиану и моду, метод k ближайших соседей и многократную импутацию цепными уравнениями. Импутер обучали только на обучающей части внутри фолдов, а на тесте пропуски заполняли по ближайшим соседям из обучающей части, без использования исхода и других тестовых наблюдений.

Рабочей выбрана импутация методом k ближайших соседей с пятью соседями. Она превосходила остальные стратегии по качеству модели и заметно не искажала форму распределений, при сравнении распределений до и после импутации критерием Колмогорова-Смирнова значимых расхождений не выявлено. Эта проверка ориентировочная, поскольку сравниваемые выборки вложены друг в друга, и трактовать ее как доказательство совпадения распределений нельзя. Число соседей сравнивали отдельно (k от 3 до 11) по out-of-fold ROC-AUC на обучающей части, лучшим оказалось k = 5.

\subsection{Статистический анализ}\label{subsec2_7}

Нормальность проверяли критерием Шапиро-Уилка при n менее 50 и критерием Колмогорова-Смирнова в варианте Лиллиефорса при n не менее 50. Количественные показатели при нормальном распределении описывали средним и стандартным отклонением с 95\% доверительным интервалом, при отличии от нормального - медианой и квартилями. Группы сравнивали t-критерием Стьюдента, U-критерием Манна-Уитни или W-критерием Бруннера-Мюнцеля в зависимости от нормальности и равенства дисперсий (критерий Левена). Доли в четырехпольных таблицах сравнивали критерием хи-квадрат Пирсона при минимальной ожидаемой частоте более 10 и точным критерием Фишера при ожидаемой частоте менее 10. Доверительный интервал для долей считали методом Клоппера-Пирсона, размер эффекта - отношением шансов с 95\% доверительным интервалом. При множественных сравнениях применяли поправку Бенджамини-Хохберга. Площади под ROC-кривыми двух моделей на одних и тех же пациентах сравнивали критерием DeLong для связанных кривых.

\subsection{Разработка моделей}\label{subsec2_8}

Выборку один раз стратифицированно разделили в отношении 80 на 20, обучающая часть 218 пациентов, тестовая 55. Разбиение выполнено до любой предобработки, с фиксированным зерном генератора случайных чисел. Всю предобработку (импутацию, масштабирование, кодирование категориальных признаков, балансировку) выполняли внутри пайплайна и обучали только на обучающих фолдах.

Выбор стратегий и семейств моделей проводили целиком на обучающей части, к отложенному тесту на этом этапе не обращались. Сравнили семь семейств (логистическая регрессия, случайный лес, LightGBM, CatBoost, XGBoost, метод опорных векторов, метод k ближайших соседей) по осям импутации, балансировки, кодирования категориальных и преобразования количественных признаков, критерий сравнения - out-of-fold ROC-AUC на обучающей части.

В дальнейшую работу взяли четыре семейства, представляющих разные классы моделей, логистическую регрессию как линейную модель и эталон сравнения для методологической гипотезы, случайный лес как метод бэггинга, XGBoost и CatBoost как варианты градиентного бустинга, второй из которых обрабатывает категориальные признаки нативно. Метод опорных векторов и метод k ближайших соседей отстали заметно (out-of-fold ROC-AUC 0.677 и 0.631 против 0.713-0.780 у остальных) и дальше не рассматривались. LightGBM исключен как избыточный по отношению к XGBoost, обе модели реализуют градиентный бустинг над деревьями, а различие между ними (0.724 против 0.716) мало на фоне разброса оценок по фолдам.

Дисбаланс классов учитывали взвешиванием классов, синтетические методы (SMOTE, SMOTENC, BorderlineSMOTE, ADASYN) сравнивали отдельно. Гиперпараметры подбирали байесовским поиском во вложенной кросс-валидации, качество процедуры оценивали по внешней пятифолдовой кросс-валидации, а внутри каждого обучающего фолда гиперпараметры подбирались по отдельной трехфолдовой кросс-валидации.

\subsection{Оценка качества}\label{subsec2_9}

Основные метрики выбраны независимыми от точки отсечения, ROC-AUC и PR-AUC с 95\% доверительными интервалами по бутстрэпу (2000 повторов). Именно они отвечают на вопрос, насколько клинический профиль разделяет группы, и не зависят от произвольного выбора порога.

Калибровку вероятностей проверяли по out-of-fold калибровочным кривым и индексу Бриера, сырые вероятности сравнивали с методами Платта и изотонической регрессии, правило применения задали заранее, калибровку используем только там, где она снижает индекс Бриера не менее чем на 0.005.

Дополнительно, для наглядности структуры ошибок, приводим матрицу ошибок. Точку отсечения выбирали нейтральным критерием, по максимуму индекса Юдена на out-of-fold вероятностях обучающей части, затем фиксировали и применяли к отложенному тесту, который использовали один раз. Такая точка не служит клинической рекомендацией, поскольку исход устанавливается ретроспективно, задачи скрининга здесь не возникает, и отсечение служит только для иллюстрации соотношения ошибок первого рода (ложноположительные) и второго рода (ложноотрицательные).

\subsection{Интерпретация}\label{subsec2_10}

Вклад признаков оценивали методом SHAP для всех четырех моделей, LinearSHAP для логистической регрессии, TreeSHAP для случайного леса и XGBoost, встроенный расчет значений SHAP для CatBoost. Вклады категорий, закодированных методом one-hot, суммировали обратно к исходному признаку, важность признака измеряли средним абсолютным значением SHAP.

Важность признаков разбирали на двух уровнях, отвечающих на разные вопросы.

Первый уровень - согласие моделей между собой. Мы проверяли, совпадает ли ранжирование признаков у разных моделей или оно зависит от выбора алгоритма. Оценивали попарной ранговой корреляцией Спирмена между полными ранжированиями, а также пересечением топов, для чего при порогах топ-5, топ-10 и топ-15 считали, сколько признаков попадают в топ по SHAP одновременно у всех четырех моделей.

Зависимость этого пересечения от выборки проверяли бутстрэпом, 200 повторов, в каждом обучающая часть пересэмплировалась с возвращением, все четыре модели переобучались с зафиксированными гиперпараметрами, SHAP рассчитывался на не попавших в выборку наблюдениях.

Второй уровень - сопоставление машинного обучения с классической статистикой. Мы проверяли, совпадают ли признаки, общие для топ-10 всех моделей, со списком показателей, значимых при сравнении групп. Оба списка получены независимо, на состав признаков, поданных моделям, результаты статистического сравнения не влияли. Дополнительно для каждого показателя из Таблицы~\ref{tab1} считали, в скольких из четырех моделей он попал в топ-10 по SHAP.

\subsection{Программное обеспечение и отчетность}\label{subsec2_11}

Анализ выполнен на Python 3.11 с фиксированными версиями библиотек, код открыт. Отчетность построена по рекомендациям STROBE для наблюдательных исследований, для части, касающейся разработки моделей, дополнительно учтены рекомендации TRIPOD+AI.

\section{Результаты}\label{sec3}

\subsection{Характеристика участников}\label{subsec3_1}

Из исходной когорты в 289 пациентов исключены 16 без установленного исхода. Аналитическая выборка составила 273 пациента, 87 с рецидивирующим течением (31.9\%) и 186 с единичным эпизодом. Стратифицированное разбиение дало обучающую часть из 218 пациентов (69 с рецидивирующим течением, 149 с единичным эпизодом) и тестовую из 55 (18 и 37 соответственно).

В Таблице~\ref{tab1} собраны показатели со значимым различием между группами после поправки Бенджамини-Хохберга, всего 10.

\begin{table}[h]
\caption{Признаки со значимым различием между группами (поправка Бенджамини-Хохберга, q менее 0.05)}\label{tab1}%
\begin{tabular}{@{}lllll@{}}
\toprule
Показатель & Единичный (n=186) & Рецидивирующий (n=87) & q & ОШ (95\% ДИ)\\
\midrule
Суточная доза инсулина, ед. & 34 [0; 45] & 44 [35; 54] & $<$0.0001 & -\\
Без инсулинотерапии, \% & 28.9 & 1.1 & $<$0.0001 & -\\
ЛПВП, ммоль/л & 1.09 [0.83; 1.38] & 1.26 [1.06; 1.58] & 0.0065 & -\\
Длительность диабета, лет & 6 [0; 12.5] & 9 [4; 15] & 0.0139 & -\\
HbA1c, \% & 12.56 [10.28; 13.64] & 10.90 [8.91; 12.66] & 0.0196 & -\\
Без ретинопатии, \% & 65.4 & 48.1 & 0.0196 & -\\
Без поражения почек (С0), \% & 56.3 & 25.8 & 0.0196 & -\\
Без альбуминурии (А0), \% & 52.7 & 27.8 & 0.0196 & -\\
Прием алкоголя, \% & 15.3 & 29.1 & 0.0234 & 2.27 [1.23; 4.20]\\
Полинейропатия, \% & 52.1 & 69.0 & 0.0281 & 2.05 [1.17; 3.57]\\
\botrule
\end{tabular}
\footnotetext{Количественные показатели даны как медиана [Q1; Q3]. Для категориальных приведена доля указанного уровня от числа известных значений в группе, пропуски в знаменатель не включены. Отношение шансов рассчитано только для бинарных показателей.}
\end{table}

По виду инсулинотерапии главное различие определяется не долей помповой терапии, а долей пациентов без инсулинотерапии, 28.9\% в группе единичного эпизода против 1.1\% при рецидивирующем течении. Это отражает структурный сигнал дебюта заболевания, поскольку первый эпизод по определению не может быть отнесен к рецидивирующему течению.

\subsection{Клиническая гипотеза: прием алкоголя}\label{subsec3_2}

Значение признака известно у 269 пациентов из 273. Среди пациентов с рецидивирующим течением алкоголь за сутки до эпизода принимали 29.1\% (25 из 86), в группе единичного эпизода 15.3\% (28 из 183). Минимальная ожидаемая частота в четырехпольной таблице равна 16.9, поэтому по протоколу применен критерий хи-квадрат Пирсона, статистика 7.01 при одной степени свободы, p = 0.008. Нулевую гипотезу отклоняем. Отношение шансов составило 2.27 (95\% ДИ 1.23-4.20), то есть прием алкоголя накануне повышает шансы рецидивирующего течения примерно вдвое. После поправки Бенджамини-Хохберга связь сохраняет значимость (q = 0.023). Направление эффекта согласуется с данными литературы \cite{ref7,ref9,ref16}.

\subsection{Разделяющая способность}\label{subsec3_3}

\begin{table}[h]
\caption{Разделяющая способность по вложенной кросс-валидации}\label{tab2}%
\begin{tabular}{@{}lll@{}}
\toprule
Модель & ROC-AUC & SD\\
\midrule
Случайный лес & 0.775 & 0.019\\
XGBoost & 0.748 & 0.038\\
CatBoost & 0.735 & 0.050\\
Логистическая регрессия & 0.713 & 0.029\\
\botrule
\end{tabular}
\footnotetext{SD - стандартное отклонение по внешним фолдам вложенной кросс-валидации.}
\end{table}

\begin{table}[h]
\caption{Разделяющая способность на отложенном тесте}\label{tab3}%
\begin{tabular}{@{}ll@{}}
\toprule
Модель & ROC-AUC (95\% ДИ)\\
\midrule
Логистическая регрессия & 0.802 [0.672; 0.912]\\
XGBoost & 0.769 [0.630; 0.888]\\
CatBoost & 0.739 [0.587; 0.868]\\
Случайный лес & 0.724 [0.576; 0.851]\\
\botrule
\end{tabular}
\footnotetext{Доверительные интервалы получены бутстрэпом (2000 повторов).}
\end{table}

Оценки на тесте лежат в диапазоне 0.72-0.80, по вложенной кросс-валидации - в диапазоне 0.71-0.78. Доверительные интервалы на тесте широкие (в среднем 0.26 по ширине) и сильно перекрываются, поэтому ранжировать модели только по результату на тесте нельзя. Ранжирование к тому же разошлось с вложенной кросс-валидацией на противоположное, случайный лес, лидер на кросс-валидации, на тесте оказался последним, а логистическая регрессия, последняя на кросс-валидации, показала на тесте наилучший результат. При 55 пациентах в тестовой части такое расхождение ожидаемо и само по себе служит предостережением против выбора модели по одной отложенной выборке.

PR-AUC на обучающей части держится в диапазоне 0.52-0.55 при базовой линии случайного классификатора около 0.32.

Признаком переобучения мы считали систематическое превышение оценок, полученных на обучающей части, над оценкой на отложенном тесте за пределами 95\% доверительных интервалов. Такого превышения нет, у всех четырех моделей оценка вложенной кросс-валидации попадает внутрь доверительного интервала тестовой оценки. Дополнительно у случайного леса согласуются две схемы ресемплинга на обучающей части, out-of-bag оценка 0.749 [0.683; 0.815] и out-of-fold оценка 0.772 [0.710; 0.833] взаимно накрываются доверительными интервалами. Обе проверки не доказывают отсутствия переобучения, а показывают, что расхождения между оценками не выходят за пределы статистической неопределенности при данном объеме выборки.

Попарное сравнение четырех моделей критерием DeLong с поправкой Бенджамини-Хохберга на 6 сравнений не выявило ни одной значимой пары ни на out-of-fold обучающей части, ни на отложенном тесте. Наименьшие значения получены при сравнении логистической регрессии со сложными моделями на обучающей части (p после поправки 0.178). Это не доказывает равенства моделей, поскольку при таком объеме выборки мощность критерия мала, но и не дает оснований выделить какую-либо модель как лучшую.

\subsection{Методологическая гипотеза}\label{subsec3_4}

Заранее заданную методологическую гипотезу проверяли отдельно (Таблица~\ref{tab4}). На обучающей части все три сложные модели номинально выше логистической регрессии (разница 0.04-0.05), но после поправки различие незначимо ни у одной. На отложенном тесте знак разницы обратный, логистическая регрессия выше всех трех сложных моделей на 0.03-0.08, и здесь различия тоже незначимы.

\begin{table}[h]
\caption{Сложные модели против логистической регрессии, критерий DeLong}\label{tab4}%
\begin{tabular}{@{}lllll@{}}
\toprule
Выборка & Сложная модель & AUC сложной & AUC логрега & p (BH)\\
\midrule
OOF & Случайный лес & 0.772 & 0.719 & 0.09\\
OOF & XGBoost & 0.761 & 0.719 & 0.09\\
OOF & CatBoost & 0.764 & 0.719 & 0.09\\
Тест & Случайный лес & 0.724 & 0.802 & 0.50\\
Тест & XGBoost & 0.769 & 0.802 & 0.55\\
Тест & CatBoost & 0.739 & 0.802 & 0.50\\
\botrule
\end{tabular}
\footnotetext{OOF - out-of-fold предсказания на обучающей части. Поправка Бенджамини-Хохберга применена отдельно внутри каждой выборки, по три сравнения.}
\end{table}

Нулевую гипотезу отклонить не удалось, превосходство сложных моделей над логистической регрессией на этой выборке не продемонстрировано. Подчеркнем корректную трактовку, это не доказательство эквивалентности моделей, поскольку при 218 пациентах на обучающей части и 55 на тесте мощность критерия ограничена.

\subsection{Калибровка и структура ошибок}\label{subsec3_5}

Сырые вероятности оказались откалиброваны лучше, чем любой из вариантов пересчета. Индекс Бриера на out-of-fold предсказаниях составил 0.178-0.189 для сырых вероятностей, 0.185-0.192 после метода Платта и 0.190-0.193 после изотонической регрессии. Заданное заранее правило (применять калибровку только при снижении индекса Бриера не менее чем на 0.005) не выполнилось ни для одной модели, поэтому во всем дальнейшем анализе используются сырые вероятности.

Отметим побочный эффект, калибровка снижала и разделяющую способность (у случайного леса с 0.772 до 0.744 после метода Платта). Само по себе сигмоидное преобразование монотонно и порядок вероятностей, а с ним и ROC-AUC, изменить не может. Снижение дает не оно, а реализация процедуры, калибровочная обертка заново обучает базовую модель на фолдах кросс-валидации и усредняет их предсказания, поэтому меняется сама модель, а не только шкала вероятностей. При выборке такого объема эта потеря заметна.

Структура ошибок на отложенном тесте в точке отсечения по индексу Юдена приведена в Таблице~\ref{tab5}.

\begin{table}[h]
\caption{Матрица ошибок на отложенном тесте, точка отсечения по индексу Юдена}\label{tab5}%
\begin{tabular}{@{}llllllll@{}}
\toprule
Модель & Порог & TP & FP & FN & TN & Чувств. & Специф.\\
\midrule
Логистическая регрессия & 0.315 & 16 & 13 & 2 & 24 & 0.889 & 0.649\\
CatBoost & 0.250 & 14 & 15 & 4 & 22 & 0.778 & 0.595\\
XGBoost & 0.380 & 13 & 10 & 5 & 27 & 0.722 & 0.730\\
Случайный лес & 0.415 & 11 & 10 & 7 & 27 & 0.611 & 0.730\\
\botrule
\end{tabular}
\footnotetext{Тестовая часть, 18 пациентов с рецидивирующим течением и 37 с единичным эпизодом. FP - ошибка первого рода, FN - ошибка второго рода.}
\end{table}

Матрица показывает то, что не видно по ROC-AUC, при близкой разделяющей способности модели по-разному распределяют ошибки между двумя типами. Логистическая регрессия пропускает лишь 2 случая рецидивирующего течения из 18, но ошибочно относит к нему 13 пациентов из 37. Случайный лес ведет себя противоположно, ложных тревог меньше, зато пропущено 7 случаев из 18. XGBoost дает наиболее сбалансированную картину. Разброс подчеркивает нестабильность точки отсечения при таком объеме выборки, поэтому содержательные выводы мы строим на метриках, не зависящих от порога.

\subsection{Отрицательные результаты}\label{subsec3_6}

Влияние преобразований количественных признаков (ручное клиническое преобразование и преобразование Йео-Джонсона против только стандартизации) проверено для всех четырех моделей. Выигрыша нет ни у одной, у логистической регрессии расхождение между вариантами 0.004 по ROC-AUC, у CatBoost 0.005, у XGBoost результат совпадает до третьего знака во всех трех вариантах (0.712). Последнее ожидаемо, поскольку деревья решений инвариантны к монотонным преобразованиям признаков и логарифм с преобразованием Йео-Джонсона не меняют порядок значений, а значит и точки разбиения. Единственное заметное расхождение дал случайный лес (0.773 без преобразования, 0.756 при ручном, 0.775 при Йео-Джонсоне), но и оно к преобразованию отношения не имеет, вариант ручного преобразования собирается из трех блоков признаков и меняет порядок колонок, а случайный лес при фиксированном зерне отбирает случайные подмножества признаков по их индексам.

Стекинг с логистической регрессией в качестве мета-модели над out-of-fold предсказаниями базовых моделей дал лучшую комбинацию (случайный лес и XGBoost) с ROC-AUC 0.7699, что не превышает лучшую одиночную модель (0.772). Добавление базовых моделей качество не повышало.

Синтетическая балансировка также не дала прироста. ROC-AUC у синтетических методов составила 0.715-0.725 (SMOTENC 0.716, BorderlineSMOTE 0.715, ADASYN 0.722, SMOTE поверх one-hot 0.725), тогда как без балансировки 0.729, а при взвешивании классов 0.737. Индекс Бриера при этом растет с 0.187 до 0.193-0.200, то есть калибровка ухудшается. Лучший результат дал несинтетический прием взвешивания классов, он и применялся во всех финальных моделях.

\subsection{Значимость признаков}\label{subsec3_7}

Полные ранжирования признаков по среднему абсолютному значению SHAP согласуются между моделями умеренно (Таблица~\ref{tab6}). Ожидаемо выше всего согласие между случайным лесом и XGBoost, обе модели строят ансамбли деревьев. Ниже всего согласие у CatBoost с остальными.

\begin{table}[h]
\caption{Попарная ранговая корреляция Спирмена между ранжированиями SHAP}\label{tab6}%
\begin{tabular}{@{}lllll@{}}
\toprule
 & CatBoost & Логрег & Лес & XGBoost\\
\midrule
CatBoost & 1.00 & 0.49 & 0.58 & 0.47\\
Логрег & 0.49 & 1.00 & 0.71 & 0.51\\
Лес & 0.58 & 0.71 & 1.00 & 0.91\\
XGBoost & 0.47 & 0.51 & 0.91 & 1.00\\
\botrule
\end{tabular}
\end{table}

Умеренное согласие полных ранжирований объясняется концами списков, где абсолютные значения SHAP малы и ранги неустойчивы. Верхняя часть списков совпадает заметно сильнее, причем размер совпадения зависит от порога отсечения. В топ-5 у всех четырех моделей общими оказываются два признака, вид инсулинотерапии и суточная доза инсулина. В топ-10 таких признаков становится пять, добавляются HbA1c, ЛПВП и глюкоза при поступлении. В топ-15 их восемь. Выделенного числа здесь нет, и сам рост совпадения при ослаблении порога ожидаем, признаков всего 25, и при пороге во весь список общими оказались бы все. Информативен поэтому не рост, а состав совпадения при жестком пороге, где случайное пересечение маловероятно.

Совпадение при топ-10 держится на выборке неравномерно. При бутстрэпе обучающей части (200 повторов) вид инсулинотерапии и суточная доза инсулина остаются общими для всех моделей в 92\% и 89\% повторов, тогда как HbA1c, ЛПВП и глюкоза при поступлении - в 32-40\%. У отдельных моделей эти три признака попадают в топ-10 существенно чаще, в 47-72\% повторов. Разрыв между двумя величинами отражает свойство самого критерия, требование одновременного присутствия у всех четырех моделей чувствительнее к возмущению данных, чем ранжирование каждой модели по отдельности. Та же пара признаков, что устойчива к пересэмплированию, совпадает у всех моделей уже при более жестком пороге топ-5. Без пересэмплирования картина повторяется, при расчете SHAP на отложенном тесте вместо обучающей части ранжирования отдельных моделей почти не меняются (коэффициент Спирмена 0.97-1.00), а совпадение при топ-10 сокращается с пяти признаков до четырех. Сказывается и объем выборки, при 218 наблюдениях исключение нескольких пациентов заметно меняет ранжирование.

Различия между моделями приходятся на признаки за пределами этого совпадения и отражают их устройство. Логистическая регрессия дополнительно поднимает электролиты (натрий и калий), случайный лес и XGBoost - длительность диабета и общий холестерин, CatBoost - тип диабета.

Сопоставление с классической статистикой показывает, что из пяти признаков, общих для топ-10 всех моделей, четыре значимы и при сравнении групп - вид инсулинотерапии, суточная доза инсулина, HbA1c и ЛПВП. Пятый, глюкоза при поступлении, значимых различий между группами при одномерном сравнении не показал, но в топ-10 вошел у всех четырех моделей.

Полная картина соответствия приведена в Таблице~\ref{tab7}. Показатели с наименьшими значениями q входят в топ-10 у всех моделей, а по мере роста q число моделей убывает. Согласие двух подходов сильнее всего для самых значимых показателей и постепенно ослабевает.

\begin{table}[h]
\caption{Соответствие между значимостью при сравнении групп и вкладом по SHAP}\label{tab7}%
\begin{tabular}{@{}lll@{}}
\toprule
Показатель & q & В топ-10 SHAP, моделей из 4\\
\midrule
Вид инсулинотерапии & $<$0.0001 & 4\\
Суточная доза инсулина & $<$0.0001 & 4\\
ЛПВП & 0.0065 & 4\\
HbA1c & 0.0196 & 4\\
Глюкоза при поступлении & не значима & 4\\
Поражение почек (ХБП, С) & 0.0196 & 3\\
Мочевина при поступлении & не значима & 3\\
Общий холестерин & не значим & 3\\
Длительность диабета & 0.0139 & 2\\
Ретинопатия & 0.0196 & 2\\
Альбуминурия (ХБП, А) & 0.0196 & 2\\
Полинейропатия & 0.0281 & 1\\
Прием алкоголя & 0.0234 & 0\\
\botrule
\end{tabular}
\footnotetext{Приведены показатели, значимые при сравнении групп, и показатели, вошедшие в топ-10 по SHAP не менее чем у трех моделей. Значение q взято из Таблицы~\ref{tab1}.}
\end{table}

Расхождения в нижней части таблицы ожидаемы, поскольку две величины измеряют разное. Критерий сравнения групп оценивает изолированное различие по одному показателю, а средний абсолютный SHAP - вклад признака в предсказание, усредненный по всей когорте и рассчитанный при наличии остальных 24 переменных. Бинарный признак с умеренной распространенностью получает при таком усреднении малое значение, даже если для своей подгруппы пациентов он информативен.

\section{Обсуждение}\label{sec4}

Мы оценили, насколько по клиническому профилю пациента при поступлении можно отличить рецидивирующее течение ДКА от единичного эпизода, и выявили связанные с исходом факторы. Разделяющая способность держится в диапазоне 0.71-0.78 на вложенной кросс-валидации и 0.72-0.80 на отложенном тесте.

Один из главных результатов - общие признаки у двух независимых подходов, многомерного анализа SHAP и одномерного статистического сравнения. Вид инсулинотерапии, суточная доза инсулина, HbA1c, ЛПВП и глюкоза при поступлении входят в топ-10 по SHAP у всех четырех моделей машинного обучения, то есть не зависят от выбора алгоритма. Четыре из этих пяти признаков значимы и при сравнении групп по статистическому протоколу. Это важно методологически, классическая статистика оценивает каждый признак изолированно, а SHAP - его вклад внутри многомерной модели с учетом остальных переменных, поэтому подходы опираются на разные допущения.

Однако подходы совпадают не во всем, и эти различия носят закономерный характер. Глюкоза при поступлении вошла в общий для моделей топ-10, хотя при одномерном сравнении групп значимой не была, то есть многомерная модель выделила показатель, который сам по себе групп не разделяет. Есть и обратные случаи, часть значимых показателей попадает в топ-10 лишь у некоторых моделей (Таблица~\ref{tab7}). Статистический критерий и SHAP измеряют разное, поэтому расхождение между ними несет дополнительную информацию, а не указывает на ошибку. Именно такое перекрестное сравнение двух независимых методологий составляет методическую сторону работы и на малой выборке помогает отделить общий для методов сигнал от особенностей каждого из них.

Клиническая гипотеза подтвердилась, прием алкоголя за сутки до эпизода связан с рецидивирующим течением, отношение шансов 2.27 (95\% ДИ 1.23-4.20). Это переносит известный из литературы фактор риска на российскую популяцию \cite{ref7,ref9,ref16}. Отметим ограничение трактовки, дизайн поперечный, поэтому речь идет о связи, а не о доказанной причинности.

Методологическая гипотеза не подтвердилась. Практический вывод, при выборке такого объема усложнение модели не окупается, и логистическая регрессия предпочтительна как самая простая и интерпретируемая, с естественной калибровкой вероятностей. Случайный лес и XGBoost дают то же качество за счет нелинейных взаимодействий, CatBoost удобен нативной обработкой категориальных признаков, но выигрыша ни один из них не приносит. Показательно и то, что рейтинг моделей на отложенном тесте оказался обратным рейтингу на вложенной кросс-валидации, при 55 пациентах в тестовой части различия между моделями неотличимы от случайных колебаний. Этот отрицательный результат согласуется с систематическим обзором \cite{ref17}, который на материале других клинических задач также не обнаружил преимущества машинного обучения над логистической регрессией.

Прямое сравнение разделяющей способности с опубликованными моделями некорректно, поскольку они решают другую задачу. На когорте взрослых с диабетом 1 типа получена ROC-AUC около 0.82 \cite{ref11}, до 0.85 на когорте детей \cite{ref10}, но обе работы предсказывают первый эпизод или госпитализацию в заданном временном окне на крупных зарубежных базах.

\section{Ограничения}\label{sec5}

Главное ограничение - дизайн. Исход установлен ретроспективно, поэтому на момент измерения предикторов он в большинстве случаев уже реализовался. Модель диагностическая, она не предсказывает будущее событие, а разделяет уже сложившиеся группы по клиническому профилю.

Длительность диабета различалась между группами, медиана 9 лет [4; 15] при рецидивирующем течении против 6 лет [0; 12.5] при единичном эпизоде (q = 0.0139). Это различие может отражать не риск, а само определение исхода, чем дольше длится заболевание, тем за больший период мог быть зарегистрирован повторный эпизод, а общего для всех пациентов интервала наблюдения в данных нет. Построение прогностической модели требует проспективной когорты с датами эпизодов, временем наблюдения и цензурированием, что составляет естественное продолжение работы.

Совпадение признаков у всех моделей подтверждено лишь частично. При бутстрэпе обучающей части вид инсулинотерапии и суточная доза инсулина остаются общими для всех моделей в 92\% и 89\% повторов, а HbA1c, ЛПВП и глюкоза при поступлении - в 32-40\%, поэтому последние три следует считать требующими подтверждения. Полученные частоты представляют собой верхнюю границу, гиперпараметры при пересэмплировании фиксировались, и изменчивость их подбора в оценку не заложена. Перевыборка из одной когорты, кроме того, ничего не говорит о воспроизводимости на данных другого центра.

Отдельно отметим структурный сигнал в признаке вида инсулинотерапии. Доля пациентов без инсулинотерапии составила 28.9\% в группе единичного эпизода против 1.1\% при рецидивирующем течении, в нашей когорте отсутствие инсулинотерапии на момент госпитализации почти всегда соответствует дебюту заболевания, а первый эпизод по определению не может быть отнесен к рецидивирующему течению. Вклад этого признака поэтому частично воспроизводит определение исхода, и как самостоятельный фактор риска он трактоваться не может, хотя именно он чаще остальных признаков оказывается общим для всех моделей.

Остальные ограничения связаны с данными. Выборка одноцентровая и малая, внешней валидации на независимой популяции нет. Доля пропусков в части лабораторных показателей высокая, до 58\%, и хотя импутация заметно не исказила распределения, неопределенность сохраняется. Ряд признаков, названных в литературе ведущими для рецидивирующего течения (психиатрические расстройства, приверженность терапии, социально-экономические обстоятельства), в наших данных отсутствует, и это ограничивает достижимую разделяющую способность.

\section{Выводы}\label{sec6}

Клинический профиль пациента при поступлении разделяет рецидивирующее и единичное течение ДКА с умеренной способностью, ROC-AUC 0.71-0.78 по вложенной кросс-валидации и 0.72-0.80 на отложенном тесте. В топ-10 по SHAP одновременно у всех четырех моделей входят пять признаков, вид инсулинотерапии, суточная доза инсулина, HbA1c, ЛПВП и глюкоза при поступлении. Четыре из этих пяти значимы и при сравнении групп, то есть выделяются двумя независимыми методологическими путями. При бутстрэпе обучающей части общими для всех моделей стабильно остаются два из них, вид инсулинотерапии и суточная доза инсулина (92\% и 89\% повторов), остальные три - в 32-40\%.

Клиническая гипотеза подтвердилась, прием алкоголя за сутки до эпизода связан с рецидивирующим течением, отношение шансов 2.27 (95\% ДИ 1.23-4.20). Методологическая гипотеза не подтвердилась, превосходство сложных моделей над логистической регрессией продемонстрировать не удалось, стекинг и синтетическая балансировка преимущества также не дали, что при данном объеме данных говорит в пользу простых интерпретируемых моделей.

Модель диагностическая, а не прогностическая на временном горизонте. Для прогностической модели нужны проспективная когорта с датами эпизодов и временем наблюдения, а также внешняя валидация на данных другого центра.

\backmatter

\bmhead{Дополнительные материалы}

Полные таблицы результатов (характеристика когорты, сравнение стратегий предобработки, подбор гиперпараметров, калибровка, точки отсечения и матрицы ошибок, попарные сравнения DeLong, стекинг, синтетическая балансировка, частоты, с которыми признаки оказываются общими для всех моделей при бутстрэпе) и рисунки SHAP доступны в открытом репозитории проекта.

\bmhead{Благодарности}

Авторы благодарят врачей, предоставивших обезличенные клинические данные.

\section*{Declarations}

\bmhead{Финансирование}

Исследование выполнено без внешнего финансирования.

\bmhead{Конфликт интересов}

Авторы заявляют об отсутствии конфликта интересов, относящегося к содержанию настоящей статьи.

\bmhead{Этическое одобрение и согласие на участие}

[Заполняется после получения решения этического комитета. Требуемая форма: полное название комитета, номер протокола и дата, а также указание на соответствие Хельсинкской декларации 1964 года с последующими поправками. Для ретроспективного исследования на обезличенных данных допустима формулировка об освобождении от необходимости одобрения с теми же реквизитами.]

\bmhead{Согласие на публикацию}

Не применимо, работа выполнена на обезличенных данных, индивидуальные сведения о пациентах не публикуются.

\bmhead{Доступность данных}

Данные пациентов не могут быть опубликованы, поскольку это нарушило бы конфиденциальность. Доступ возможен по обоснованному запросу к автору для переписки в рамках этических ограничений.

\bmhead{Доступность материалов}

Не применимо.

\bmhead{Доступность кода}

Код анализа открыт и доступен по адресу \url{https://github.com/ovelsad/dka-recurrence-ml}, включает фиксированные версии библиотек и инструкции по воспроизведению.

\bmhead{Вклад авторов}

О. Садыхов: концепция исследования, подготовка и анализ данных, статистический анализ, разработка и валидация моделей, интерпретация результатов, подготовка первого варианта рукописи. В. Леоненко: концепция и дизайн исследования, научное руководство, методология анализа, проверка результатов, критический пересмотр и редактирование рукописи. Оба автора прочитали и одобрили окончательный вариант.

\begin{thebibliography}{17}

\bibitem{ref1} Benoit, S.R., Zhang, Y., Geiss, L.S., Gregg, E.W., Albright, A.: Trends in diabetic ketoacidosis hospitalizations and in-hospital mortality - United States, 2000-2014. MMWR Morb. Mortal. Wkly. Rep. \textbf{67}(12), 362-365 (2018). \url{https://doi.org/10.15585/mmwr.mm6712a3}

\bibitem{ref2} Desai, D., Mehta, D., Mathias, P., Menon, G., Schubart, U.K.: Health care utilization and burden of diabetic ketoacidosis in the U.S. over the past decade: a nationwide analysis. Diabetes Care \textbf{41}(8), 1631-1638 (2018). \url{https://doi.org/10.2337/dc17-1379}

\bibitem{ref3} Santos, S.S., Ramaldes, L.A.L., Dualib, P.M., Gabbay, M.A.L., Sa, J.R., Dib, S.A.: Increased risk of death following recurrent ketoacidosis admissions: a Brazilian cohort study of young adults with type 1 diabetes. Diabetol. Metab. Syndr. \textbf{15}(1), 85 (2023). \url{https://doi.org/10.1186/s13098-023-01054-5}

\bibitem{ref4} Gibb, F.W., Teoh, W.L., Graham, J., Lockman, K.A.: Risk of death following admission to a UK hospital with diabetic ketoacidosis. Diabetologia \textbf{59}(10), 2082-2087 (2016). \url{https://doi.org/10.1007/s00125-016-4034-0}

\bibitem{ref5} Mohler, R., Lotharius, K., Moothedan, E., Goguen, J., Bandi, R., Beaton, R., Knecht, M., Mejia, M.C., Khoury, M., Sacca, L.: Factors contributing to diabetic ketoacidosis readmission in hospital settings in the United States: a scoping review. J. Diabetes Complications \textbf{38}(10), 108835 (2024). \url{https://doi.org/10.1016/j.jdiacomp.2024.108835}

\bibitem{ref6} Dhatariya, K.K., Glaser, N.S., Codner, E., Umpierrez, G.E.: Diabetic ketoacidosis. Nat. Rev. Dis. Primers \textbf{6}(1), 40 (2020). \url{https://doi.org/10.1038/s41572-020-0165-1}

\bibitem{ref7} Brandstaetter, E., Bartal, C., Sagy, I., Jotkowitz, A., Barski, L.: Recurrent diabetic ketoacidosis. Arch. Endocrinol. Metab. \textbf{63}(5), 531-535 (2019). \url{https://doi.org/10.20945/2359-3997000000158}

\bibitem{ref8} Kanzwl, S.M.O.S., Alhajri, A.H.M., Mohammed, Y.J.A., et al.: A comparative effectiveness of intravenous fluids and insulin regimens in the acute management of diabetic ketoacidosis (DKA) and hypoglycemia: a systematic review. Cureus \textbf{17}(10), e94902 (2025). \url{https://doi.org/10.7759/cureus.94902}

\bibitem{ref9} Cooper, H., Tekiteki, A., Khanolkar, M., Braatvedt, G.: Risk factors for recurrent admissions with diabetic ketoacidosis: a case-control observational study. Diabet. Med. \textbf{33}(4), 523-528 (2016). \url{https://doi.org/10.1111/dme.13004}

\bibitem{ref10} Williams, D.D., Ferro, D., Mullaney, C., Skrabonja, L., Barnes, M.S., Patton, S.R., Lockee, B., Tallon, E.M., Vandervelden, C.A., Schweisberger, C., Mehta, S., McDonough, R., Lind, M., D'Avolio, L., Clements, M.A.: An ``all-data-on-hand'' deep learning model to predict hospitalization for diabetic ketoacidosis in youth with type 1 diabetes: development and validation study. JMIR Diabetes \textbf{8}, e47592 (2023). \url{https://doi.org/10.2196/47592}

\bibitem{ref11} Li, L., Lee, C.C., Zhou, F.L., Molony, C., Doder, Z., Zalmover, E., Sharma, K., Juhaeri, J., Wu, C.: Performance assessment of different machine learning approaches in predicting diabetic ketoacidosis in adults with type 1 diabetes using electronic health records data. Pharmacoepidemiol. Drug Saf. \textbf{30}(5), 610-618 (2021). \url{https://doi.org/10.1002/pds.5199}

\bibitem{ref12} Rudin, C.: Stop explaining black box machine learning models for high stakes decisions and use interpretable models instead. Nat. Mach. Intell. \textbf{1}(5), 206-215 (2019). \url{https://doi.org/10.1038/s42256-019-0048-x}

\bibitem{ref13} Stiglic, G., Kocbek, P., Fijacko, N., Zitnik, M., Verbert, K., Cilar, L.: Interpretability of machine learning-based prediction models in healthcare. WIREs Data Min. Knowl. Discov. \textbf{10}(5), e1379 (2020). \url{https://doi.org/10.1002/widm.1379}

\bibitem{ref14} Ahmed, S., Kaiser, M.S., Hossain, M.S., Andersson, K.: A comparative analysis of LIME and SHAP interpreters with explainable ML-based diabetes predictions. IEEE Access \textbf{13}, 37370-37388 (2025). \url{https://doi.org/10.1109/ACCESS.2024.3422319}

\bibitem{ref15} Emi-Johnson, O.G., Nkrumah, K.J.: Predicting 30-day hospital readmission in patients with diabetes using machine learning on electronic health record data. Cureus \textbf{17}(4), e82437 (2025). \url{https://doi.org/10.7759/cureus.82437}

\bibitem{ref16} Bradford, A.L., Crider, C.C., Xu, X., Naqvi, S.H.: Predictors of recurrent hospital admission for patients presenting with diabetic ketoacidosis and hyperglycemic hyperosmolar state. J. Clin. Med. Res. \textbf{9}(1), 35-39 (2017). \url{https://doi.org/10.14740/jocmr2792w}

\bibitem{ref17} Christodoulou, E., Ma, J., Collins, G.S., Steyerberg, E.W., Verbakel, J.Y., Van Calster, B.: A systematic review shows no performance benefit of machine learning over logistic regression for clinical prediction models. J. Clin. Epidemiol. \textbf{110}, 12-22 (2019). \url{https://doi.org/10.1016/j.jclinepi.2019.02.004}

\end{thebibliography}

\end{document}
